\subsection{ভিন্ন সারণিতে SQL অনুশীলন}
প্রতিটি টেবিল স্বাধীন নমুনা। নির্দেশ লিখে নিজের ফল যাচাই করো।
\subsubsection{Employee}\begin{flushleft}\small
\begin{tabular}{|l|l|l|l|}
\hline

\chSixHeader
ID & Name & Designation & Salary\\
\hline
101 & Anika & Manager & 40000\\\hline
102 & Shafin & Officer & 25000\\\hline
103 & Fahim & Accountant & 32000\\\hline
104 & Kamal & Officer & 48000\\\hline
105 & Javed & Officer & 52000\\\hline
106 & Marjana & Officer & 29000\\
\hline

\end{tabular}
\end{flushleft}
\begin{enumerate}
\item Officer পদে Salary 30000-এর বেশি কিন্তু 55000-এর সমান বা কম এমন কর্মীদের তথ্য Name অনুযায়ী সাজাও।
\end{enumerate}
\subsubsection{Student}\begin{flushleft}\small
\begin{tabular}{|l|l|l|l|}
\hline

\chSixHeader
Name & Class & Roll & Section\\
\hline
Mizan & 9 & 1 & morning\\\hline
David & 9 & 2 & morning\\\hline
Parimal & 8 & 2 & day\\\hline
Rubel & 8 & 1 & day\\\hline
Sumi & 9 & 4 & morning\\\hline
Tania & 10 & 5 & morning\\
\hline

\end{tabular}
\end{flushleft}
\begin{enumerate}
\item Class 9 এবং morning Section-এর শিক্ষার্থীদের দেখাও।
\item Class 9 অথবা 10-এর শিক্ষার্থীদের দেখাও।
\end{enumerate}
\subsubsection{Patient}\begin{flushleft}\small
\begin{tabular}{|l|l|l|l|l|}
\hline

\chSixHeader
Name & Gender & Age & Blood & Contact\\
\hline
Zahid & Male & 45 & O+ & 02918205\\\hline
Sabina & Female & 50 & B+ & 02918305\\\hline
Malek & Male & 60 & AB+ & 02918504\\\hline
Rahim & Male & 35 & O+ & 02918702\\\hline
Anika & Female & 25 & O+ & 02918901\\
\hline

\end{tabular}
\end{flushleft}
\begin{enumerate}
\item Male রোগীদের মধ্যে Age কমপক্ষে 50 হলে Name, Age, Blood ও Contact দেখাও।
\item Anika-এর Gender ভুলভাবে Male রাখা হলে Female দিয়ে সংশোধন করো। একই নামের একাধিক রোগী থাকলে কী দিয়ে রেকর্ড শনাক্ত করবে?
\end{enumerate}
Worker ও Product-এর পৃথক অনুশীলন অধ্যায়ের শেষের SQL অনুশীলনে দেওয়া হয়েছে।
